i\documentclass[instructions.tex]{subfiles}
\begin{document}

\chapter{Est-il avantageux de vivre long-temps ?}

\primo Vous voudriez vivre long-temps : mais si vous aviez un vrai désir de posséder Dieu, vous soupireriez, avec saint Paul, après l'heureux moment qui doit vous unir à Jésus-Christ\footnote{Desiderium habens dissolvi, et esse cum Christo. Philip. 1.}. Loin d'appréhender la mort, vous gémiriez, comme David, de la longue durée de votre vie\footnote{Heu mihi! quia incolatus meus prolongatus est. Ps. 119.}.

Vous voudriez vivre long-temps, et à peine avez-vous pensé à vivre. De quoi sert-il de vivre, si l'on ne vit que pour pécher ? Vous voudriez vivre long-temps ; mais on peut dire, en général, qu'une mort avancée semble être préférable à une longue vie. Elle est en un sens plus avantageuse aux méchants et aux gens de bien : avantageuse aux méchants, qui multiplient leurs crimes en vivant long-temps. Ils sont bien plus malheureux, dit saint Ambroise, de vivre dans le péché et d'y persévérer, que d'être surpris de la mort après l'avoir commis\footnote{Gravius est ad peccandum vivere, quam in peccato mori. S. Ambr., lib. de bon. mort.}.

\secundo Une mort avancée paraît avantageuse aux gens de bien, parce qu'ils jouiront plus tôt de Dieu ; ce qui a fait dire à saint Jean Chrysostome, qu'on doit se réjouir à la mort d'un enfant et à la mort d'un homme craignant Dieu, parce qu'ils sont délivrés des misères de cette vie, et vont jouir d'un bonheur éternel.

Des gens de bien, en vivant plus long-temps, peuvent, à la vérité, donner plus de gloire au Créateur, et s'acquitter de plus en plus de leurs dettes envers Dieu ; mais ils peuvent aussi en contracter de nouvelles, et ils ne donneraient jamais tant de gloire à Dieu par leurs bonnes œuvres, qu'ils lui en raviraient par les moindres péchés. S'il est avantageux de vivre dans la sainteté, il est encore plus avantageux de mourir saintement. Une sainte mort met notre salut hors de danger : une sainte vie nous y dispose, mais elle n'assure pas encore entièrement le salut.

Quelque sainte que soit notre vie, sommes-nous assurés de persévérer ? Qui peut savoir toutes les occasions qui sont capables d'ébranler notre courage ? Qui peut se flatter de n'être pas entraîné par le mauvais exemple ? Qui peut prévoir le temps et le lieu qui nous cache une Dalila qui pourra nous séduire et nous perdre ? Il y a toujours à craindre, tandis qu'un vaisseau est en mer : souvent la tempête le brise dans l'endroit où le calme l'avait laissé en sûreté le même jour.

La vie présente est une mer pleine d'écueils, un labyrinthe sujet à mille égarements, une terre qui dévore ses habitants.

Ainsi, demander une longue vie, c'est demander une longue suite de dangers, de tentations et de combats\footnote{Longam tentationem petit qui longam vitam petit. Curé de Chartreux de Trèves.}.

Combien de personnes, hélas ! à qui une longue vie a été l'occasion de leur perte ! Combien d'hommes vertueux, qui semblaient n'avoir rien à craindre, se sont pervertis tout-à-coup et ont péri ! \emph{Ne faites rien pour m'empêcher de mourir}, disait au lit de la mort le bienheureux Louis de Gonzague, \emph{car je ne sais ce qui peut m'arriver, si je vis plus long-temps.}

Concluez de toutes ces réflexions : que vous devez conserver votre vie sans inquiétudes, et avoir plus d'empressement de vivre dans la sainteté, que de vivre long-temps ; que vous devez vous détacher de la vie, et ne point tant craindre la mort. Votre sort est entre les mains de Dieu : soumettez-vous avec confiance à sa volonté, qui a fixé le nombre de vos jours, et profitez du temps qu'il vous donne.

Je me soumets à Dieu pour la longueur ou la brièveté de ma vie. Mais je lui demanderai tous les jours la grâce de bien vivre. Et je me défierai constamment de mon extrême faiblesse.

Ô mon Dieu ! affermissez-moi dans le bien, et faites-moi la grâce de terminer ma course mortelle d'une manière sainte et précieuse à vos yeux.

\end{document}
